\chapter{Resumen}

\begin{abstract}
Las subestaciones electricas de distribucion emiten de forma continua medidas
de tension, corriente, frecuencia, temperatura y calidad de potencia a traves
de sistemas SCADA y PMU. Sin embargo, las fallas son eventos infrecuentes y
las etiquetas historicas son escasas, lo que dificulta entrenar detectores
supervisados. Ademas, los datos crudos no pueden salir del recinto por
restricciones de ciberseguridad, ancho de banda y propiedad. Se necesita por
lo tanto un sistema de deteccion \emph{temprana, no supervisada} y compatible
con una arquitectura federada.

Este trabajo propone un pipeline de cuatro componentes:
\begin{enumerate}[leftmargin=*,itemsep=2pt]
  \item Un \textbf{gemelo digital ligero} que predice el comportamiento
        electrico esperado de la subestacion a partir de la topologia y
        las consignas de operacion, generando ademas datos sinteticos
        etiquetados para entrenamiento y validacion.
  \item Un \textbf{detector no supervisado} basado en autoencoder LSTM sobre
        el residuo multivariado observado-esperado, mas \emph{conformal
        prediction} para fijar el umbral de alarma con una \emph{tasa de
        falsa alarma garantizada}.
  \item Una arquitectura de \textbf{aprendizaje federado} (FedAvg) entre
        tres subestaciones virtuales (prototipo; en despliegue se usan
        cinco o mas), preservando los datos en origen y
        cumpliendo con las restricciones regulatorias.
  \item Una capa de \textbf{priorizacion post-evento} para ordenar cuadrillas
        tras un evento climatico extremo (vientos, calor, frente de lluvia),
        reduciendo SAIDI y la exposicion a multas regulatorias.
\end{enumerate}

Los resultados del prototipo ejecutable muestran que el conformal
prediction cumple la tasa de falsa alarma garantizada: FPR empirico
de $4.86\%$ cuando el objetivo conformal es $\alpha = 0.05$, con
latencia mediana de deteccion de $0$~s. Con configuracion reducida
(modelo de 32 unidades, ventana de 15 muestras, 2 dias de
simulacion), el detector alcanza un recall de $76.3\%$
(centralizado) y $80.5\%$ (federado), con F1 de $0.80$ y $0.87$
respectivamente: dentro del rango que la literatura reporta para
la configuracion completa. La version federada no cede
detectabilidad frente a la centralizada y mueve $0$~MB de datos
crudos fuera de las subestaciones.

El pipeline es publicable como software vendible a distribuidoras
(Chilquinta, CGE, Saesa) y transmisoras (Engie Transelec), con un
caso de uso directo en la deteccion de eventos climaticos extremos
que la SEC y el SERNAC han comenzado a fiscalizar con cargos directos
a las concesionarias tras el frente de Valparaiso de julio de 2026.
\end{abstract}

\noindent\textbf{Palabras clave:} gemelo digital, subestacion electrica,
SCADA, PMU, deteccion de anomalias, autoencoder, conformal prediction,
aprendizaje federado, SAIDI, ciberseguridad electrica.

\begin{abstract}[english]
\foreignlanguage{english}{\textbf{Digital twin, distribution substation,
unsupervised early warning, SCADA/PMU, conformal prediction, federated
learning.}}

We propose a four-component pipeline for early unsupervised fault detection
in distribution substations: a lightweight digital twin of the substation
that produces synthetic labeled data; an LSTM autoencoder trained on the
observed-vs-expected multivariate residual, with a conformal-prediction
threshold that yields a guaranteed false-alarm rate; a three-Site FedAvg
federated training loop that keeps raw data on premise and complies with
NERC CIP / CNE cybersecurity regulation; and a post-event prioritization
layer to dispatch field crews after extreme weather events. On a two-day
simulation with eight fault types representative of the Chilean grid, the
detector achieves an empirical FPR of $4.86\%$ against a $\alpha = 0.05$
conformal target, a recall of $76.3\%$ (centralized) and $80.5\%$
(federated) with F1 of $0.80$ and $0.87$, and a median detection latency
of $0$~s. The federated model matches the centralized one (AUC $0.906$ vs
$0.878$) while moving zero raw data off premise. The pipeline is
productizable as a SaaS for distribution utilities and transmission
owners.
\end{abstract}